% TOP_PROBLEM: {"problemId": "problem.grigorchuk-gap-conjecture-on-growth-of-groups", "canonicalTitle": "Grigorchuk Gap Conjecture", "releaseRank": 282, "primaryDomain": "algebra_representation_category", "primaryDomainLabel": "Algebra, representation & category theory", "displayStatus": "open_with_solved_subcases", "releaseStatus": "open_with_solved_subcases"}
% !TeX program = lualatex
% Definition and sources only; this file is not an LLM solution attempt.
\documentclass[11pt]{article}
\usepackage[margin=25mm]{geometry}
\usepackage{fontspec}
\setmainfont{Latin Modern Roman}
\usepackage{amsmath,amssymb,mathtools}
\usepackage{unicode-math}
\setmathfont{Latin Modern Math}
\usepackage{xurl,hyperref}
\hypersetup{colorlinks=true,urlcolor=blue}
\setcounter{secnumdepth}{0}
\setlength{\parindent}{0pt}
\setlength{\parskip}{5pt}
\setlength{\emergencystretch}{3em}
\begin{document}
\section*{\texorpdfstring{Grigorchuk Gap Conjecture}{Grigorchuk Gap Conjecture}}\label{282}
\subsection{Definitions and mathematical statement}
For a finite generating set $X$ of $G$, let $\gamma_{G,X}(n)$ count elements represented by words of length at most $n$ in $X\cup X^{-1}$. Write $f\preceq g$ if an integer $C\ge1$ gives $f(n)\le Cg(Cn)$ for every $n\ge1$, and $f\prec g$ when $f\preceq g$ but not $g\preceq f$. The conjecture is
\[
 \gamma_{G,X}\prec\exp(\sqrt n)\quad\Longrightarrow\quad
 \gamma_{G,X}(n)=O(n^d)\text{ for some }d.
\]
Equivalently the conclusion is virtual nilpotence: a finite-index subgroup is nilpotent. Strict comparison of growth classes is essential; an ordinary pointwise little-o comparison can hold between functions in the same growth class after rescaling the argument.
\par\smallskip\textit{Formulation and concept references:} \href{https://arxiv.org/pdf/1202.6044}{[S1]}.

\subsection{Short English statement}
Is there a genuine gap between polynomial group growth and the growth class of the exponential of the square root?
\subsection{Sources}
\begin{itemize}
\item[S1]Rostislav Grigorchuk, On the Gap Conjecture Concerning Group Growth (2012/2014), Conjecture 1; Eberhard and Maini (2025), Section 1 comparison convention.
\url{https://arxiv.org/pdf/1202.6044}.
\textit{Formulation source.}
\end{itemize}
\end{document}
